\section{¿Qué es Dota 2 y qué es farmear?}

\begin{frame}{Dota 2 en dos minutos}
  \begin{itemize}
    \item Videojuego de estrategia por equipos: dos equipos de cinco héroes.
    \item Objetivo: destruir la base enemiga.
    \item Para volverse más fuerte, el héroe necesita \textbf{oro}.
    \item El oro se obtiene principalmente matando \textbf{creeps}:
      \begin{itemize}
        \item \textbf{Creeps de línea:} aparecen cada 30 segundos.
        \item \textbf{Creeps neutrales:} aparecen cada 60 segundos en la jungla.
      \end{itemize}
  \end{itemize}
  \vspace{0.5em}
  \begin{block}{Analogía}
    Es como un tablero donde las monedas aparecen solas cada cierto tiempo en caminos fijos y en rincones. El jugador que recoja más monedas en menos tiempo gana ventaja.
  \end{block}
\end{frame}

\begin{frame}{¿Por qué el farmeo se puede estudiar con matemáticas?}
  \begin{itemize}
    \item Las apariciones siguen reglas temporales \textbf{deterministas} y documentadas.
    \item Los valores de oro por creep son \textbf{conocidos}.
    \item El tiempo de desplazamiento entre campamentos es \textbf{medible}.
  \end{itemize}
  \vspace{0.8em}
  \begin{alertblock}{Consecuencia}
    El farmeo no es azar: es un \textbf{problema de optimización} con reglas claras. Por eso se puede modelar, calcular y mejorar.
  \end{alertblock}
\end{frame}